\section{Actual line fibers and model norms at genuine proper centers}

\begin{lemma}[Actual valuative model frame and actual original generic line]
\label{lem:scheme-line-actual-center-frame}
\lean{EnvelopeProperty.Berkovich.RelativeSchemeValuedPoint.genericLineModule}
\lean{EnvelopeProperty.Berkovich.RelativeSchemeValuedPoint.centerLineModule}
\lean{EnvelopeProperty.Berkovich.RelativeSchemeValuedPoint.centerGenericMap}
\lean{EnvelopeProperty.Berkovich.RelativeSchemeValuedPoint.centerLineGenericIso}
\lean{EnvelopeProperty.Berkovich.RelativeSchemeValuedPoint.exists_centerLineUnitIso}
\lean{EnvelopeProperty.Berkovich.RelativeSchemeValuedPoint.centerFrameGenericIso}
\leanok
\uses{lem:actual-local-spectrum-line-triviality}
Let $y$ be an original valued point of a proper scheme over bounded
original coefficients, let $V_y$ be its actual residue-field valuation
ring and $c_y:\operatorname{Spec}V_y\to X$ the genuine proper center
lift. For every original invertible module sheaf $\mathcal L$, the
actual line $c_y^*\mathcal L$ has a global original frame. Pulling
that frame along $\operatorname{Spec}\kappa(y)\to\operatorname{Spec}V_y$
gives an actual global frame of the original generic line
$y^*\mathcal L$, using the genuine valuative triangle and native
pullback comparisons.
\end{lemma}
\begin{proof}
Original invertibility is preserved by actual module pullback, so
$c_y^*\mathcal L$ is invertible. The valuation ring $V_y$ is local;
apply local-spectrum line triviality. The actual valuative generic
triangle is $j_yc_y=y$. Thus the native pullback composition followed
by its native equality comparison identifies $j_y^*c_y^*\mathcal L$
with $y^*\mathcal L$. Pull back the actual global valuative frame and
use the actual structure-unit pullback comparison. Composing with
the inverse of this actual generic comparison is the required frame
on the entire original generic line.
\end{proof}

\begin{definition}[The genuine frame norm on the actual generic line]
\label{def:scheme-line-actual-center-frame-norm}
\lean{EnvelopeProperty.Berkovich.RelativeSchemeValuedPoint.centerFrameCoefficient}
\lean{EnvelopeProperty.Berkovich.RelativeSchemeValuedPoint.centerFrameLineNorm}
\lean{EnvelopeProperty.Berkovich.RelativeSchemeValuedPoint.centerFrameLineNorm_zero}
\lean{EnvelopeProperty.Berkovich.RelativeSchemeValuedPoint.centerFrameLineNorm_add_le_max}
\lean{EnvelopeProperty.Berkovich.RelativeSchemeValuedPoint.centerFrameLineNorm_eq_zero_iff}
\lean{EnvelopeProperty.Berkovich.RelativeSchemeValuedPoint.centerFrameLineNorm_smul}
\leanok
\uses{lem:scheme-line-actual-center-frame}
For an actual valuative frame $e:c_y^*\mathcal L\simeq\mathcal O$,
let $e_\eta:y^*\mathcal L\simeq\mathcal O$ be the actual generic
frame just constructed. On the entire original generic-line global
section module, set
\[
 \lvert s\rvert_e=\lvert\Gamma(e_\eta)(s)\rvert_y,
\]
using the actual canonical identification
$\Gamma(\operatorname{Spec}\kappa(y),\mathcal O)=\kappa(y)$.
This vanishes exactly for the zero section, satisfies the
nonarchimedean triangle inequality, and scales by the original
residue-field valuation under original scalar multiplication.
\end{definition}
\begin{proof}
Both the actual generic frame and the actual affine structure-ring
identification are isomorphisms, so the coefficient is zero exactly
when $s=0$. Their additive and scalar identities transport the
original residue-field valuation's zero, multiplication and
nonarchimedean inequalities to the stated section module. This uses
the original fiber and actual pullback sheaves, not an auxiliary
one-dimensional vector space.
\end{proof}
